\documentclass[conference]{IEEEtran}

\usepackage{amsmath}
\usepackage{booktabs}
\usepackage{graphicx}
\usepackage{array}
\usepackage{xcolor}
\usepackage{url}

\definecolor{pendingred}{RGB}{170,20,20}
\newcommand{\resultpending}{\textcolor{pendingred}{\textbf{[FINAL RESULT]}}}
\newcommand{\ppoints}{percentage points}

\title{A Two-Stage Plant Disease Recognition Assistant:\
Localization, Classification, Robustness, and Cross-Domain Evaluation}

\author{
\IEEEauthorblockN{Federico Volpi}
\IEEEauthorblockA{
M.Sc. in Computer Engineering\\
University of Pisa, Pisa, Italy
}
\and
\IEEEauthorblockN{Rafay Zia}
\IEEEauthorblockA{
M.Sc. in Computer Engineering\\
University of Pisa, Pisa, Italy
}}

\begin{document}

\maketitle

\begin{abstract}
Plant disease recognition in field photographs is difficult because symptoms can occupy only a
small image region and acquisition conditions differ substantially from controlled datasets. This
work presents an intelligent assistant that first localizes symptomatic regions and then classifies
the corresponding plant--disease condition. The system combines a class-agnostic YOLO11n lesion
detector with one of three classifiers: a convolutional neural network trained from scratch,
EfficientNetB0, or MobileNetV3-Large. A disease-aware YOLO11n provides a single-stage comparison.
PlantSeg is the primary field-image dataset, while PlantVillage supplies a taxonomy-matched
controlled-domain benchmark; PlantDoc is retained as optional backup data. Model stability is
assessed by five-fold stratified cross-validation over the PlantSeg training pool, and final systems
are evaluated on the untouched official test split. The study additionally measures synthetic
corruption robustness and Grad-CAM overlap with lesion masks. The best two-stage system reaches
59.02\% macro F1 on the official holdout. Five-fold cross-validation confirms this ranking with a
mean macro F1 of 59.39\% ($\pm$1.99 \ppoints), showing the result is stable across different
train/validation partitions rather than an artifact of one split. Severe Gaussian blur is the
strongest tested corruption for the classifier pipelines, while JPEG compression most strongly
affects standalone YOLO. Grad-CAM obtains 24.2\% mean top-quantile IoU with the annotated lesions.
On the matched 21-class PlantVillage benchmark, the strongest classifier pipeline reaches
39.03\% macro F1, a decrease of 25.72 \ppoints{} from the equivalent PlantSeg subset. These
experiments evaluate accuracy together with localization, stability,
explanation alignment, and domain sensitivity rather than relying on a single aggregate score.
\end{abstract}

\begin{IEEEkeywords}
plant disease recognition, object detection, image classification, YOLO, transfer learning,
cross-validation, domain shift, explainable artificial intelligence, robustness
\end{IEEEkeywords}

\section{Introduction}

Plant diseases threaten crop quality and yield, and timely diagnosis often depends on specialist
knowledge that may not always be immediately available. Image-based recognition can support an
initial screening step using ordinary cameras. However, field photographs contain clutter,
variable lighting, scale changes, occlusion, and multiple plant regions. A model evaluated only on
centered leaves and uniform backgrounds may therefore appear reliable without transferring to
natural scenes. This gap has also been observed in prior plant-disease research: models attaining
very high accuracy on PlantVillage suffered a substantial decrease on images acquired under
different conditions~\cite{mohanty2016plantvillage}.

This project investigates whether explicit symptom localization supports more robust field-image
classification. The proposed assistant separates \emph{where the disease evidence is} from
\emph{which disease it represents}. A class-agnostic detector first identifies a suspected lesion;
a classifier then assigns one of 115 PlantSeg plant--disease labels to the resulting crop. The
assistant also reports confidence and a Grad-CAM explanation. It is intended as decision support,
not as a substitute for laboratory testing, a plant pathologist, or treatment advice.

The study addresses five research questions:

\begin{itemize}
    \item \textbf{RQ1:} How does the two-stage localization--classification pipeline compare with
    a single disease-aware YOLO detector on PlantSeg field images?
    \item \textbf{RQ2:} Which classifier provides the best and most stable macro-level performance
    across the official protocol and five cross-validation folds?
    \item \textbf{RQ3:} How does performance change when the PlantSeg-trained pipeline is applied
    to a taxonomy-matched PlantVillage subset?
    \item \textbf{RQ4:} Which common image corruptions cause the greatest degradation, and do the
    classifier architectures degrade differently?
    \item \textbf{RQ5:} To what extent does classifier attention overlap the lesion masks supplied
    by PlantSeg?
\end{itemize}

The contributions are a reproducible two-stage pipeline for lesion localization and 115-class
disease recognition; a controlled comparison with disease-aware YOLO11n; complete five-fold
retraining to estimate split variability; a 21-class PlantSeg--PlantVillage comparison; and
synthetic-corruption and mask-based explanation analyses. Together, these experiments cover the
data selection, cleaning, transformation, modeling, evaluation, and prototype stages of the
Knowledge Discovery in Databases (KDD) process.

\section{Related Work}

\subsection{Plant-Disease Image Classification}

PlantVillage made large-scale supervised plant-disease classification accessible by collecting
more than 54,000 healthy and diseased leaf images. Mohanty \emph{et al.} reported 99.35\% accuracy
over 38 crop--disease classes using convolutional networks~\cite{mohanty2016plantvillage}. The
dataset is nevertheless controlled: leaves are centered, image dimensions are uniform, and
background variation is limited. Their evaluation on differently acquired images fell to slightly
above 31\%, illustrating that within-dataset performance is not sufficient evidence for field use.

PlantDoc was later introduced to represent non-laboratory imagery, including background noise,
varying illumination, multiple leaves, and lower-resolution symptoms~\cite{singh2020plantdoc}.
PlantSeg further expands in-the-wild supervision by pairing diverse plant-disease images with
pixel-level lesion masks across 115 disease categories and 34 plant hosts~\cite{wei2024plantseg}.
These resources motivate treating controlled and natural images as distinct domains rather than
merging them without an explicit experimental role.

Transfer learning is useful when individual disease classes contain few images. EfficientNet uses
compound scaling to balance network depth, width, and input resolution~\cite{tan2019efficientnet}.
MobileNetV3 combines hardware-aware architecture search with efficient mobile building blocks,
targeting a favorable accuracy--latency trade-off~\cite{howard2019mobilenetv3}. This work compares
both ImageNet-initialized architectures with a small CNN trained from scratch on identical lesion
crops.

\subsection{Localization and Visual Explanation}

Modern one-stage detectors make localization practical in an interactive pipeline. YOLO11 supports
object detection with compact variants such as YOLO11n and provides COCO-pretrained
weights~\cite{jocher2024yolo11}. The present study uses the same nano architecture in two roles: a
single-class lesion detector and a 115-class disease-aware detector. The first delegates disease
discrimination to a specialized classifier; the second is simpler at inference but learns
localization and fine-grained classification jointly.

Grad-CAM uses class-score gradients flowing into the last convolutional layer to construct a coarse,
class-discriminative activation map without modifying or retraining the network
~\cite{selvaraju2017gradcam}. Such a map can reveal whether a classifier attends near visible
symptoms, but it is neither a causal explanation nor a segmentation prediction. PlantSeg masks
enable a quantitative attention-alignment test rather than relying solely on attractive overlays.

\subsection{Robustness and Domain Shift}

Common-corruption benchmarks demonstrate that neural classifiers may be unstable under blur,
noise, compression, and other non-adversarial changes~\cite{hendrycks2019robustness}. This project
adapts that principle to plant imagery by testing six corruption families at five severity levels.
It complements the synthetic study with a cross-dataset experiment: the unchanged PlantSeg-trained
pipeline is evaluated on PlantVillage classes that have defensible disease-level equivalents. This
design distinguishes sensitivity to controlled perturbations from a broader acquisition-domain
change.

\section{Use Cases and System Architecture}

\subsection{Use-Case Scenarios}

In the primary scenario, a user uploads an RGB photograph of a symptomatic plant. The assistant
localizes a suspected lesion, displays the crop or bounding region, predicts a supported disease,
and presents confidence and likely alternatives. A Grad-CAM overlay exposes regions associated with
the prediction, while a warning discourages reliance on low-confidence output. The result is a
preliminary screening aid; an agronomic expert remains responsible for diagnosis and intervention.

A secondary analysis scenario serves developers and researchers. Saved predictions, confusion
matrices, corruption curves, and explanation examples can be inspected through the experiment
notebook. This makes errors attributable to localization, classification, class imbalance, or
domain shift easier to separate.

The assistant is closed-set: it always selects among PlantSeg disease categories after either a
detector crop or full-image fallback. Healthy plants, unknown diseases, and images far outside the
training distribution are therefore not reliably rejected. Confidence is communicated as a model
score rather than diagnostic certainty.

\subsection{Architecture}

Fig.~\ref{fig:architecture} summarizes the two-stage inference path. The detector returns zero or
more lesion boxes. When boxes exist, their union is expanded by a 10\% margin and cropped; otherwise
the full image is passed to the classifier. The selected classifier produces a softmax distribution,
from which the top disease, confidence, and alternatives are obtained. Grad-CAM is computed for the
predicted class within the same crop and mapped back to full-image coordinates.

\begin{figure}[t]
\centering
\fbox{\begin{minipage}{0.92\columnwidth}
\centering
Uploaded RGB image\\[2pt]
$\downarrow$\\[2pt]
Class-agnostic YOLO11n lesion detector\\[2pt]
$\swarrow$ detection \hspace{12mm} no detection $\searrow$\\[-1pt]
Union box + 10\% margin \hspace{7mm} Full image fallback\\[2pt]
$\searrow$ \hspace{32mm} $\swarrow$\\[2pt]
Selected 115-class classifier\\[2pt]
$\downarrow$\\[2pt]
Disease, confidence, alternatives, and Grad-CAM\\[2pt]
$\downarrow$\\[2pt]
User interface and low-confidence warning
\end{minipage}}
\caption{Two-stage plant disease assistant. The single-stage baseline replaces the detector and
classifier with one disease-aware YOLO11n.}
\label{fig:architecture}
\end{figure}

The standalone comparison replaces the cascade with one disease-aware YOLO11n. Its
highest-confidence box provides the image-level disease prediction. If the detector returns no box,
the image is counted as incorrect; consequently, both conditional classification quality and
coverage affect its end-to-end score.

\subsection{Model Components}

\textbf{Class-agnostic YOLO11n:} COCO-pretrained YOLO11n is fine-tuned to predict one \texttt{lesion}
class. Each target box encloses the nonzero region of a binary PlantSeg lesion mask. This formulation
asks the detector to learn localization without separating 115 disease labels.

\textbf{Disease-aware YOLO11n:} the standalone model uses the same architecture, initialization,
images, and mask-derived geometry, but each box carries the image-level PlantSeg disease ID. The
comparison therefore changes the detector label task while holding the localization target fixed.

\textbf{Baseline CNN:} the from-scratch classifier contains four convolutional blocks with 32, 64,
128, and 256 channels. Every block applies a $3\times3$ convolution, batch normalization, ReLU, and
$2\times2$ max pooling. Adaptive average pooling, dropout with probability 0.3, and a 115-output
linear layer produce the logits.

\textbf{EfficientNetB0 and MobileNetV3-Large:} both networks begin from ImageNet weights
~\cite{deng2009imagenet}. Their original heads are replaced by 115-output linear layers and the
complete networks are fine-tuned on detector-generated crops. EfficientNetB0 provides a compact
compound-scaled model; MobileNetV3-Large provides a mobile-oriented alternative.

\section{Datasets and KDD Process}

\subsection{Dataset Roles}

Table~\ref{tab:datasets} summarizes the deliberately separate roles of the three datasets. They are
not merged because their taxonomies, annotation types, and visual distributions differ.

\begin{table*}[t]
\caption{Dataset characteristics and experimental roles. Counts describe the audited local
representations used by the project.}
\label{tab:datasets}
\centering
\resizebox{\textwidth}{!}{%
\begin{tabular}{lrrlll}
\toprule
Dataset & Images & Classes & Acquisition & Available supervision & Role in this work \\
\midrule
PlantSeg & 7,774 & 115 & In-the-wild & Image label and lesion mask & Primary train/validation/test, localization, Grad-CAM \\
PlantVillage & 54,305 & 38 & Controlled & Image label & Secondary domain on 21 mapped diseases \\
PlantDoc & 2,578 & 29 & In-the-wild & 8,910 boxes & Optional backup after cleaning \\
\bottomrule
\end{tabular}}
\end{table*}

\subsection{PlantSeg}

The audited PlantSeg release contains 7,774 paired RGB images and lesion masks across 34 plant types
and 115 plant--disease combinations. Its predefined split contains 5,367 training, 846 validation,
and 1,561 test images. The training split includes every class, while validation lacks class ID 68
and test lacks class ID 41. The median class contains 55 images, and class sizes range from 3 to 323.
This imbalance motivates macro F1, balanced accuracy, and per-class recall rather than accuracy alone.

The project's audit treats \texttt{Metadata.csv} as the image-label authority and PNG masks as the
localization authority. Every nonzero mask value is interpreted as lesion foreground. Fresh
enclosing boxes are derived because the supplied COCO data contain 220 image-level mask/box
disagreements, 51 out-of-bounds boxes, one image without an annotation, and incomplete category
metadata. Eight EXIF orientation cases were identified. The explanation analysis applies
orientation jointly to each image and mask before measuring overlap, while prepared YOLO views use
repair-safe copies when image encoding requires it. PlantSeg supplies all model fitting, the
official holdout, synthetic robustness images, and explanation ground truth.

\subsection{PlantVillage}

PlantVillage contains 54,305 audited color images in 38 classes, with all images at
$256\times256$ pixels. Leaves are usually centered against simple backgrounds, making it a useful
controlled benchmark but not evidence of natural-scene performance. Class sizes range from 152 to
5,507 images. The repository's official split contains 43,596 training and 10,709 test images.

The project never trains on PlantVillage. An explicit mapping retains 21 disease-level classes that
have defensible PlantSeg equivalents; healthy categories are excluded because PlantSeg has no
healthy class. The mapped official PlantVillage test view contains 6,057 images, while the same
classes yield 378 images from the PlantSeg official test split. Both views use the unchanged
PlantSeg-trained detector, classifiers, thresholds, and preprocessing. The comparison therefore
measures domain behavior rather than the benefit of additional training data.

\subsection{PlantDoc}

The audited PlantDoc representation contains 2,578 field images, 29 labels, and 8,910 bounding
boxes. It offers cluttered backgrounds, occlusion, multiple leaves, and heterogeneous resolution,
but also contains exact train/test duplicates, invalid or out-of-bounds boxes, two classes absent
from test, severe annotation imbalance, and boxes with mixed granularity. Although the original
dataset was created to study recognition outside laboratory settings~\cite{singh2020plantdoc}, this
local representation is excluded from the minimum experiment. It can be used only after grouped
deduplication, annotation repair, and explicit taxonomy mapping.

\subsection{KDD Workflow}

The project follows KDD as an iterative sequence in which findings at one stage change decisions in
the next. During \emph{domain and goal definition}, the intended use was narrowed from automated
treatment recommendation to preliminary screening of a symptomatic field photograph. This yielded
three functional outputs: a localized region, one disease prediction with alternatives and
confidence, and an explanation. It also established the safety constraint that low-confidence
results must be referred to an expert.

During \emph{data selection and understanding}, three candidate datasets were profiled for size,
class support, acquisition setting, supervision, duplicates, and annotation integrity. The audit
showed that PlantSeg is the only candidate providing both field photographs and pixel-level lesion
masks at the required taxonomy scale, so its official splits became the primary experiment.
PlantVillage was retained only as a 21-class controlled-domain test, and PlantDoc was excluded from
model fitting after its duplicate leakage, invalid boxes, mixed box semantics, and extreme class
imbalance were identified. Thus, dataset exploration directly determined each dataset's role.

During \emph{cleaning and transformation}, checksums and image decoding were verified, metadata was
used as the class authority, and every nonzero mask value was treated as lesion foreground. One
enclosing box was derived from each mask because 220 supplied boxes disagreed with masks and 51
extended beyond image bounds. The official splits were preserved, lightweight linked dataset views
were created, and only manually verified label equivalents entered the cross-domain subset. The
class-agnostic detector then generated train and validation crops with a 10\% margin; a missed
detection produced a documented full-image fallback. Classifier inputs were resized to
$224\times224$, augmented only during training, and normalized with ImageNet statistics.

During \emph{modeling and evaluation}, the same YOLO11n backbone was assigned either one lesion
class or 115 disease classes, isolating the effect of separating localization from classification.
The resulting crop was evaluated with a from-scratch CNN and two ImageNet-initialized classifiers.
Validation macro F1 selected classifier checkpoints because the audit revealed severe imbalance;
detection models used validation mAP. Five-fold reconstruction of the complete training chain
tested split stability, while the untouched official test split supplied final classification and
localization results. Cross-domain transfer, synthetic corruptions, calibration, latency, and
Grad-CAM-to-mask overlap then exposed failure modes hidden by aggregate accuracy. These findings
feed back into the next KDD iteration: rare-class sampling, blur augmentation, calibration, unknown
class rejection, and source-grouped deduplication are prioritized over simply increasing model
size. The selected lesion-YOLO and EfficientNetB0 cascade is exposed through the Streamlit upload
prototype, completing the path from discovered knowledge to an executable intelligent system.

\section{Experimental Setup}

\subsection{Split and Leakage-Control Protocol}

The main development path preserves PlantSeg's official partitions. The 5,367 training images
update parameters, the 846 validation images select checkpoints by validation performance, and the
1,561 test images remain an untouched final holdout. No test image is used for fitting, checkpoint
selection, threshold selection, or cross-validation. Evaluation is run after the training choices
are fixed.

Cross-validation is a separate stability estimate over only the official training pool. It does not
absorb or replace the official validation/test sets. This distinction permits reporting both
within-training variability and performance on the dataset publisher's holdout.

\subsection{Five-Fold Cross-Validation}

The implementation uses five-fold \texttt{StratifiedKFold} with shuffling and seed 42. Disease IDs
provide the strata. Two ultra-rare training classes cannot occur in every fold: class ID 41 has two
images and class ID 68 has four. These samples are removed from the standard stratification step and
assigned round-robin across folds. Hence, fold-level macro metrics average the classes represented in
that fold.

The complete processing chain is reconstructed independently in every fold. First, four folds form
the training partition and the fifth forms fold validation. Next, mask-derived YOLO views are built
for those partitions. New class-agnostic and disease-aware YOLO11n models are initialized from the
common COCO checkpoint and trained only on the four folds. The fold-specific lesion detector then
generates both training and validation crops. Finally, a new baseline CNN is initialized from
scratch, while new EfficientNetB0 and MobileNetV3-Large instances start from ImageNet weights; each
is trained on the four-fold crops and evaluated on held-out crops. This prevents a validation crop
from being generated by a detector that learned from the same image.

Each fold/model combination uses 50 epochs by default. Detector precision, recall, mAP50, and
mAP50--95, and classifier best validation macro F1 are aggregated as arithmetic mean, population
standard deviation, and completed-fold count. Completed fold/model outputs are reused on rerun.
The split is stratified by class, not grouped by source URL, plant instance, or perceptual duplicate
cluster; this is retained as a threat to validity.

\subsection{Training Configuration}

Table~\ref{tab:training} lists the shared training configuration. Both YOLO variants use
$640\times640$ input, batch size 16, and Ultralytics training with pretrained YOLO11n weights. The
classifiers use $224\times224$ input, batch size 32, AdamW with learning rate $3\times10^{-4}$ and
weight decay $10^{-4}$, cosine annealing, and ordinary cross-entropy. The current implementation
does not use class weighting or a balanced sampler.

Training augmentation for classifiers consists of random resized crop with scale 0.75--1.0,
horizontal flip, rotation within $\pm10^\circ$, and mild brightness, contrast, saturation, and hue
jitter. Validation uses deterministic resize and ImageNet normalization. The baseline CNN is
initialized from scratch; the other classifiers are ImageNet-initialized. Seed 42 is applied to
Python, NumPy, and PyTorch operations where supported.

\begin{table}[t]
\caption{Training configuration.}
\label{tab:training}
\centering
\resizebox{\columnwidth}{!}{%
\begin{tabular}{lll}
\toprule
Setting & YOLO variants & Classifiers \\
\midrule
Initialization & COCO YOLO11n & Scratch / ImageNet \\
Input resolution & $640\times640$ & $224\times224$ \\
Output classes & 1 or 115 & 115 \\
CV epochs & 50 per fold & 50 per fold \\
Batch size & 16 & 32 \\
Optimizer & Ultralytics automatic & AdamW \\
Learning rate & artifact-defined & $3\times10^{-4}$ \\
Weight decay & artifact-defined & $10^{-4}$ \\
Selection metric & validation mAP & validation macro F1 \\
Seed & 42 & 42 \\
\bottomrule
\end{tabular}}
\end{table}

\subsection{Evaluation Protocols}

\textbf{Primary holdout:} both YOLO models are evaluated against boxes derived from PlantSeg test
masks. The cascade passes the detector union crop to each classifier and uses full-image fallback
when no lesion exceeds the fixed confidence threshold. Standalone YOLO uses its highest-confidence
detection as the image label; no detection is incorrect.

\textbf{Cross-domain:} all three pipelines are evaluated on the 21-class PlantSeg and PlantVillage
mapped test views without PlantVillage fine-tuning. Each domain is reported separately, together
with the absolute PlantVillage-minus-PlantSeg change.

\textbf{Synthetic robustness:} Gaussian blur, brightness, contrast, JPEG compression, occlusion,
and crop/framing transformations are applied at severity levels 1--5. The clean baseline and every
corrupted condition use the same fixed PlantSeg images and pipeline.

\textbf{Grad-CAM alignment:} Grad-CAM is computed for the predicted class on the detector crop and
mapped into full-image coordinates. The resulting heatmap is compared pixel-for-pixel with the
binary PlantSeg mask. Representative successes and failures accompany aggregate metrics.

\subsection{Metrics and Reproducibility}

For class $c$, precision $P_c$ and recall $R_c$ produce
$F1_c=2P_cR_c/(P_c+R_c)$. Macro F1 is the unweighted mean over represented classes,

\begin{equation}
F1_{\mathrm{macro}}=\frac{1}{C}\sum_{c=1}^{C}F1_c,
\end{equation}

and balanced accuracy is mean per-class recall. Detection uses precision, recall, mAP50, and
mAP50--95. Coverage measures the fraction of images for which standalone YOLO emits a detection.
Expected calibration error aggregates the absolute confidence--accuracy gap over confidence bins
~\cite{guo2017calibration}. Robustness and domain changes are reported in absolute percentage
points. Explanation metrics include thresholded and top-quantile intersection-over-union (IoU),
precision, recall, mask-energy fraction, and pointing-game hit rate.

The implementation uses Python 3.13, PyTorch, torchvision, scikit-learn, and Ultralytics. Dataset
preparation, training, evaluation, and artifact generation are scripted, while commented Jupyter
notebooks expose the reasoning and results workflow. Final reporting records exact package versions,
hardware, accelerator, checkpoint paths, and run dates alongside each artifact.

\section{Results and Discussion}

Results below are populated directly from the final JSON/CSV artifacts after confirming that they
reference the selected checkpoints. Percentages are reported consistently, and fold variation is
never replaced by the best fold.

\subsection{Primary PlantSeg Comparison}

\begin{table*}[t]
\caption{End-to-end image classification on the official PlantSeg test split.}
\label{tab:primary-results}
\centering
\begin{tabular}{lccccc}
\toprule
System & Accuracy (\%) & Macro F1 (\%) & Balanced acc. (\%) & Coverage (\%) & Images/s \\
\midrule
Disease-aware YOLO11n & 61.37 & 53.32 & 49.60 & 86.48 & \textbf{28.0} \\
Lesion YOLO + baseline CNN & 41.32 & 29.04 & 29.99 & 100 & 27.2 \\
Lesion YOLO + EfficientNetB0 & 66.69 & 59.02 & 58.69 & 100 & 26.9 \\
Lesion YOLO + MobileNetV3-Large & \textbf{66.75} & 58.17 & 58.69 & 100 & 26.1 \\
\bottomrule
\end{tabular}
\end{table*}

As shown in Table~\ref{tab:primary-results}, EfficientNetB0 obtains the highest holdout macro F1,
59.02\%, while MobileNetV3-Large obtains the highest accuracy, 66.75\%. The best cascade exceeds
disease-aware YOLO macro F1 by 5.70 \ppoints. Standalone coverage is 86.48\%, so missed detections
remain one source of error even after longer training; errors among detected images account for the
remaining gap. RQ1 is therefore answered in favor of the two-stage design for classification
quality, although standalone YOLO remains operationally simpler.

\begin{table}[t]
\caption{Detection performance on the PlantSeg test split.}
\label{tab:detection-results}
\centering
\resizebox{\columnwidth}{!}{%
\begin{tabular}{lcccc}
\toprule
Detector & Precision & Recall & mAP50 & mAP50--95 \\
\midrule
Lesion YOLO11n & 85.27 & 81.94 & 88.22 & 61.52 \\
Disease-aware YOLO11n & 54.93 & 49.90 & 50.15 & 34.47 \\
\bottomrule
\end{tabular}}
\end{table}

The two detection rows represent different tasks. The class-agnostic model is correct when it
localizes a lesion, whereas the disease-aware model must also assign the correct fine-grained class.
The 38.07-point mAP50 gap is consistent with the additional fine-grained classification burden of
the disease-aware detector; it must assign the right disease as well as localize the lesion.

\subsection{Five-Fold Stability}

\begin{table*}[t]
\caption{Five-fold cross-validation over the PlantSeg official training pool. Values are
mean $\pm$ population standard deviation.}
\label{tab:cv-results}
\centering
\begin{tabular}{llcc}
\toprule
Model & Metric & Mean $\pm$ std. (\%) & Completed folds \\
\midrule
Lesion YOLO11n & mAP50 & 85.34 $\pm$ 0.56 & 5/5 \\
Lesion YOLO11n & mAP50--95 & 56.96 $\pm$ 1.00 & 5/5 \\
Lesion YOLO11n & Precision & 82.23 $\pm$ 0.70 & 5/5 \\
Lesion YOLO11n & Recall & 80.21 $\pm$ 0.70 & 5/5 \\
Disease-aware YOLO11n & mAP50 & 40.90 $\pm$ 0.66 & 5/5 \\
Disease-aware YOLO11n & mAP50--95 & 29.00 $\pm$ 0.78 & 5/5 \\
Disease-aware YOLO11n & Precision & 49.06 $\pm$ 1.62 & 5/5 \\
Disease-aware YOLO11n & Recall & 40.03 $\pm$ 2.74 & 5/5 \\
Baseline CNN & Macro F1 & 20.09 $\pm$ 0.75 & 5/5 \\
EfficientNetB0 & Macro F1 & \textbf{59.39} $\pm$ 1.99 & 5/5 \\
MobileNetV3-Large & Macro F1 & 57.27 $\pm$ 1.44 & 5/5 \\
\bottomrule
\end{tabular}
\end{table*}

Table~\ref{tab:cv-results} answers RQ2: EfficientNetB0 has the highest mean cross-validated macro
F1 (59.39\%, $\pm$1.99 \ppoints{}), the same ranking as on the official holdout, where it also leads
MobileNetV3-Large (59.02\% vs.\ 58.17\% macro F1). The standard deviations are small relative to
the means for every model, particularly the lesion detector (mAP50 varying by roughly one point
across folds), indicating that the official-split ranking is not an artifact of one favorable
train/validation partition. The five-fold classifier macro F1 values are systematically lower than
the corresponding official-holdout figures, consistent with each fold using a shorter, 50-epoch
training budget on four-fifths of the training pool rather than the longer schedule used for the
official-split models. Because two ultra-rare classes cannot appear in every fold, these estimates
should be interpreted as performance over the classes represented in each fold rather than uniform
evidence for all 115 classes.

\subsection{Cross-Domain Evaluation}

\begin{table*}[t]
\caption{Matched 21-class evaluation. Change is PlantVillage minus PlantSeg in percentage points.}
\label{tab:domain-results}
\centering
\begin{tabular}{lccccc}
\toprule
System & PlantSeg macro F1 & PlantVillage macro F1 & Change & PlantSeg ECE & PlantVillage ECE \\
\midrule
Baseline CNN & 34.52 & 10.92 & $-$23.60 & 4.89 & 14.81 \\
EfficientNetB0 & 63.63 & 35.73 & $-$27.90 & 28.86 & 42.79 \\
MobileNetV3-Large & \textbf{64.75} & \textbf{39.03} & $-$25.72 & 28.93 & 38.24 \\
Disease-aware YOLO11n & 63.93 & 17.59 & $-$46.34 & 11.75 & 12.05 \\
\bottomrule
\end{tabular}
\end{table*}

Every classifier loses substantial macro F1 when moved from PlantSeg's 378-image matched subset to
PlantVillage's 6,057-image matched subset over the same 21 shared classes, using the identical
lesion-YOLO front end and crop margin in both cases. MobileNetV3-Large remains the strongest
pipeline on PlantVillage (39.03\% macro F1) despite a 25.72-point drop. The baseline CNN loses the
fewest absolute points (23.60), but from a much lower starting accuracy; EfficientNetB0 drops the
most (27.90 points). Calibration also
degrades on the new domain: expected calibration error roughly triples for the baseline CNN (4.89\%
to 14.81\%) and increases by 14 and 9 points respectively for EfficientNetB0 and MobileNetV3-Large,
meaning confidence scores become less trustworthy exactly where accuracy is already lower. The
standalone disease-aware YOLO degrades even more sharply, with coverage itself falling from
86.8\% on PlantSeg to 43.2\% on PlantVillage (a 43.5-point drop) because PlantVillage's
studio-style, single-leaf, uncluttered images differ enough from PlantSeg's field photographs that
the detector frequently proposes no lesion box at all; its macro F1 falls by 46.34 points, the
largest shift of any system. The effect cannot be attributed to background alone because sample
frequency, symptom stage, image source, label interpretation, and detector behavior also differ
between datasets. RQ3 is therefore answered clearly: none of the tested systems generalize well
to PlantVillage's controlled photography conditions without further domain adaptation, and the
two-stage cascade's classifier accuracy degrades gracefully while the standalone detector's
coverage collapse makes it the least robust choice under domain shift.

\subsection{Synthetic-Corruption Robustness}

\begin{table*}[t]
\caption{Clean and severity-5 accuracy (\%) on the fixed 200-image PlantSeg robustness subset.
Parentheses contain the absolute change from the clean condition.}
\label{tab:robustness-results}
\centering
\resizebox{\textwidth}{!}{%
\begin{tabular}{lccccccc}
\toprule
Model & Clean & Blur & Brightness & Contrast & JPEG & Occlusion & Crop \\
\midrule
Baseline CNN & 40.0 & 13.5 (-26.5) & 20.5 (-19.5) & 25.5 (-14.5) & 28.0 (-12.0) & 36.0 (-4.0) & 36.0 (-4.0) \\
EfficientNetB0 & 65.0 & 28.5 (-36.5) & 46.5 (-18.5) & 57.5 (-7.5) & 43.5 (-21.5) & 56.5 (-8.5) & 61.5 (-3.5) \\
MobileNetV3-Large & 68.0 & 33.0 (-35.0) & 43.5 (-24.5) & 61.0 (-7.0) & 38.5 (-29.5) & 60.5 (-7.5) & 63.0 (-5.0) \\
Disease-aware YOLO11n & 57.5 & 14.0 (-43.5) & 36.0 (-21.5) & 47.0 (-10.5) & 7.5 (-50.0) & 53.5 (-4.0) & 49.0 (-8.5) \\
\bottomrule
\end{tabular}}
\end{table*}

Gaussian blur is the most damaging transformation for all three classifier pipelines. At severity
5, EfficientNetB0 falls from 65.0\% to 28.5\% accuracy and MobileNetV3-Large falls from 68.0\% to
33.0\%. JPEG compression is even more damaging to disease-aware YOLO11n, which falls from 57.5\%
to 7.5\%. Brightness and JPEG compression cause intermediate degradation for the classifier
pipelines. Crop and occlusion are comparatively well
tolerated, consistent with the random resized crop and horizontal-flip augmentation used during
classifier training. The pretrained classifiers lose more absolute accuracy to severe blur than
the baseline CNN, although they remain substantially more accurate. RQ4 therefore identifies blur
as the principal weakness of the classifier pipelines and JPEG compression as the principal
weakness of standalone YOLO.

\subsection{Grad-CAM Alignment}

\begin{table}[t]
\caption{Grad-CAM overlap with PlantSeg lesion masks.}
\label{tab:gradcam-results}
\centering
\begin{tabular}{lcc}
\toprule
Metric & Mean (\%) & Median (\%) \\
\midrule
IoU, threshold 0.5 & 21.3 & 18.9 \\
Precision, threshold 0.5 & 51.6 & 50.0 \\
Recall, threshold 0.5 & 31.3 & 28.4 \\
IoU, top 20\% & 24.2 & 22.1 \\
Precision, top 20\% & 41.2 & 34.2 \\
Recall, top 20\% & 52.7 & 48.1 \\
Mask energy fraction & 41.5 & 38.2 \\
Mask area fraction & 21.3 & 14.2 \\
Pointing-game hit rate & 58.8 & 100.0 \\
\bottomrule
\end{tabular}
\end{table}

Table~\ref{tab:gradcam-results} shows partial but meaningful spatial alignment. The
highest-activation point falls inside the annotated lesion in 58.8\% of images, while 41.5\% of
activation energy is concentrated inside masks that cover 21.3\% of the image on average. The
top-20\% heatmap reaches 24.2\% mean IoU and favors recall over precision, suggesting relatively
broad attention. RQ5 is therefore answered cautiously: classifier evidence is preferentially
concentrated near lesions, but the maps are not accurate lesion segmentations and do not establish
causal reasoning.

\subsection{Efficiency and Error Analysis}

A fair deployment comparison runs only the selected classifier in the cascade, timed end to end
(image load, YOLO inference, crop, and one classifier) at batch size 1 on the same Apple Silicon
(MPS) accelerator used for training, over 200 sampled test images after a 10-image warm-up.
Table~\ref{tab:primary-results} reports the resulting throughput. Standalone disease-aware YOLO is
the smallest model by far (2.66M parameters, 5.6~MB checkpoint) and reaches 28.0 images/s mean
throughput (27.0~ms median latency), but its 95th-percentile latency is higher (122.3~ms),
indicating occasional slow outlier images rather than uniformly fast inference. The baseline CNN
is the fastest cascade at 27.2 images/s (3.0M parameters, 10.5~MB), and
the two pretrained classifiers are close to each other and slightly slower (EfficientNetB0: 26.9
images/s, 6.7M parameters, 55.8~MB; MobileNetV3-Large: 26.1 images/s, 6.9M parameters, 58.0~MB).
The throughput ranking does not track accuracy: standalone YOLO is fastest but trails both
pretrained classifier pipelines in macro F1, while the baseline CNN is the fastest cascade and the
least accurate system. MobileNetV3-Large trades a small throughput cost for a 29-point higher macro
F1 than the baseline CNN. All four systems remain comfortably above real-time single-image
throughput requirements for an advisory or field-diagnosis tool on this hardware. The earlier
figure obtained by running all three classifiers simultaneously must not be used as a direct
latency comparison with standalone YOLO, since it measures combined rather than single-classifier
cost.

For the selected MobileNetV3-Large cascade, eight of the 114 represented test classes have zero
recall, all among the rarest classes in the training pool (support between one and nine test
images), including bell pepper frogeye leaf spot, broccoli ring spot, and peach rust. The most
frequent confusions pair visually similar symptoms on related host species rather than unrelated
diseases: squash and zucchini powdery mildew are each other's most common misclassification (13
and 11 occurrences), cucumber powdery mildew is frequently predicted as squash powdery mildew (10
occurrences), and ginger sheath blight is frequently predicted as rice sheath blight (10
occurrences). This pattern suggests the classifier has learned a symptom-level visual signature
that generalizes across closely related crops rather than a crop-specific one, which is desirable
for visual similarity but a genuine source of species-level error. Mean softmax confidence is
95.43\% for correct predictions versus 79.69\% for incorrect ones, a 15.74-point gap indicating the
model is comparatively, though not perfectly, under-confident on its errors; this partial
separation supports using a confidence threshold to flag low-certainty predictions for expert
review rather than treating every output as equally trustworthy. These observations motivate the
limitations and future work below.

\section{Prototype and Limitations}

\subsection{Prototype}

The prototype layer accepts a common RGB image format, loads the selected detector and classifier,
and follows the pipeline in Fig.~\ref{fig:architecture}. It displays the lesion crop, predicted
disease, model confidence, alternatives, and an explanation overlay. A low-confidence message
instructs the user to seek expert assessment rather than treating the result as a diagnosis. Models
are loaded once per application session so that subsequent requests incur only inference and
visualization latency. The interface demonstrates the intelligent-system workflow while preserving
the distinction between a course prototype and production agronomic software.

\subsection{Limitations and Threats to Validity}

PlantSeg is severely imbalanced, and its rarest classes cannot be represented in every official
partition or cross-validation fold. The five-fold split is class-stratified rather than grouped by
source or perceptual duplicate cluster, leaving possible source dependence. Mask-derived enclosing
boxes can be coarse when lesions are extensive or disconnected. The 21-class cross-domain mapping
covers only a subset of either taxonomy, and the domain sample counts and acquisition protocols are
unequal. Consequently, the domain comparison measures the combined effect of several dataset
differences rather than a single isolated factor.

The classifier is closed-set and has no explicit unknown or healthy rejection mechanism. Synthetic
corruptions are useful controlled probes but do not span all field conditions. Grad-CAM is coarse
and non-causal. Runtime depends on hardware, batch size, library version, and whether all classifiers
or only the selected model are executed. PlantDoc is omitted from the core experiment because its
local annotations and split require further repair. Finally, neither agronomic outcomes nor the
consequences of incorrect advice are evaluated; the system must not prescribe treatment.

\section{Conclusion}

This work designs and evaluates an intelligent assistant for recognizing plant disease in field
images. The system combines class-agnostic YOLO11n localization with EfficientNetB0 and is compared
with a disease-aware single-stage YOLO11n. On the official PlantSeg test split, the selected
cascade obtains 59.02\% macro F1, 5.70 \ppoints{} above standalone YOLO. Five-fold cross-validation
confirms this ranking is stable across different train/validation partitions (EfficientNetB0 mean
macro F1 59.39\%, $\pm$1.99 \ppoints, ahead of MobileNetV3-Large at 57.27\%, $\pm$1.44 \ppoints,
and the from-scratch baseline CNN at 20.09\%, $\pm$0.75 \ppoints). The matched PlantVillage
cross-domain evaluation shows this ranking is preserved but every system degrades substantially
under domain shift: MobileNetV3-Large remains the strongest classifier pipeline (39.03\% macro F1,
down 25.72 \ppoints{} from PlantSeg) while the standalone YOLO's usable coverage collapses from
86.8\% to 43.2\%. Severe blur is the largest corruption weakness for the classifier pipelines,
while JPEG compression causes the largest loss for standalone YOLO. Grad-CAM places 41.5\% of its
activation energy inside lesions occupying 21.3\% of the image on average.

Overall, the available results support explicit lesion localization as a useful basis for
field-image disease recognition, and five-fold cross-validation confirms the model ranking is not
an artifact of a single data split; the domain-shift results show this generalization is fragile
and the two-stage cascade degrades more gracefully than the standalone detector. Future
work should prioritize source-grouped
deduplication, better rare-class learning, calibrated unknown-class rejection, domain-diverse field
data, and direct segmentation. These extensions would strengthen both scientific validity and the
safety of an eventual deployment.

\begin{thebibliography}{00}

\bibitem{mohanty2016plantvillage}
S. P. Mohanty, D. P. Hughes, and M. Salath\'{e}, ``Using deep learning for image-based plant
disease detection,'' \emph{Frontiers in Plant Science}, vol. 7, Art. no. 1419, 2016,
doi: 10.3389/fpls.2016.01419.

\bibitem{singh2020plantdoc}
D. Singh, N. Jain, P. Jain, P. Kayal, S. Kumawat, and N. Batra, ``PlantDoc: A dataset for visual
plant disease detection,'' in \emph{Proc. 7th ACM IKDD CoDS and 25th COMAD}, 2020, pp. 249--253,
doi: 10.1145/3371158.3371196.

\bibitem{wei2024plantseg}
T. Wei, Z. Chen, X. Yu, S. Chapman, P. Melloy, and Z. Huang, ``PlantSeg: A large-scale in-the-wild
dataset for plant disease segmentation,'' arXiv:2409.04038, 2024. Dataset release:
doi: 10.5281/zenodo.17719108.

\bibitem{tan2019efficientnet}
M. Tan and Q. V. Le, ``EfficientNet: Rethinking model scaling for convolutional neural networks,''
in \emph{Proc. 36th Int. Conf. Machine Learning}, 2019, pp. 6105--6114.

\bibitem{howard2019mobilenetv3}
A. Howard \emph{et al.}, ``Searching for MobileNetV3,'' in \emph{Proc. IEEE/CVF Int. Conf.
Computer Vision}, 2019, pp. 1314--1324, doi: 10.1109/ICCV.2019.00140.

\bibitem{jocher2024yolo11}
G. Jocher and J. Qiu, ``Ultralytics YOLO11,'' version 11.0.0, 2024. [Online]. Available:
\url{https://github.com/ultralytics/ultralytics}

\bibitem{selvaraju2017gradcam}
R. R. Selvaraju, M. Cogswell, A. Das, R. Vedantam, D. Parikh, and D. Batra, ``Grad-CAM: Visual
explanations from deep networks via gradient-based localization,'' in \emph{Proc. IEEE Int. Conf.
Computer Vision}, 2017, pp. 618--626, doi: 10.1109/ICCV.2017.74.

\bibitem{hendrycks2019robustness}
D. Hendrycks and T. Dietterich, ``Benchmarking neural network robustness to common corruptions and
perturbations,'' in \emph{Proc. Int. Conf. Learning Representations}, 2019.

\bibitem{deng2009imagenet}
J. Deng, W. Dong, R. Socher, L.-J. Li, K. Li, and L. Fei-Fei, ``ImageNet: A large-scale
hierarchical image database,'' in \emph{Proc. IEEE Conf. Computer Vision and Pattern Recognition},
2009, pp. 248--255.

\bibitem{guo2017calibration}
C. Guo, G. Pleiss, Y. Sun, and K. Q. Weinberger, ``On calibration of modern neural networks,'' in
\emph{Proc. 34th Int. Conf. Machine Learning}, 2017, pp. 1321--1330.

\end{thebibliography}

\end{document}
